\chapter{La promesa que nadie te va a hacer}

\epigraph{«Lo que Dios ha unido, que no lo separe el hombre.»}{Mateo 19,6}

Antes de entrar en materia, quiero hacer una pausa y contarte algo que a mí me sorprendió cuando lo descubrí. Porque los tres capítulos que vienen ahora ---este y los dos siguientes--- giran en torno a tres palabras que suelen producir urticaria: indisolubilidad, fidelidad y fecundidad. Tres palabras que, cuando alguien las escucha juntas, tiende a pensar en reglas, prohibiciones y corsés. En cosas que la Iglesia exige y que la vida moderna ha dejado atrás.

Y sin embargo.

Si abres el Código Civil ---no el Derecho Canónico, no el catecismo, sino el Código Civil, el que regula el matrimonio para todos, creyentes o no---, te encuentras con algo llamativo. El artículo 68 dice, textualmente, que los cónyuges están obligados a vivir juntos, guardarse fidelidad y socorrerse mutuamente.\footnote{Art. 68 del Código Civil español: «Los cónyuges están obligados a vivir juntos, guardarse fidelidad y socorrerse mutuamente. Deberán, además, compartir las responsabilidades domésticas y el cuidado y atención de ascendientes y descendientes y otras personas dependientes a su cargo.»} Y el artículo 67 añade que deben respetarse y actuar en interés de la familia.\footnote{Art. 67 del Código Civil español: «Los cónyuges deben respetarse y ayudarse mutuamente y actuar en interés de la familia.»}

Fidelidad. Vida en común. Socorro mutuo. Interés de la familia. ¿Te suenan?

Son, con otras palabras, las mismas columnas sobre las que se levanta el matrimonio cristiano. No porque la Iglesia haya impuesto sus ideas al Estado ---que es lo que muchos piensan---, sino por algo mucho más interesante: porque estos principios son tan humanos, están tan inscritos en la naturaleza del amor, que el derecho civil los ha reconocido por su propio peso. La ley no necesitó consultar ningún catecismo para llegar a estas conclusiones. Bastó con observar qué es lo que hace que el amor funcione, qué es lo que necesita una familia para sostenerse, qué es lo que dos personas se deben la una a la otra cuando deciden compartir la vida.

Piénsalo un momento. Si algo tan «religioso» como la fidelidad aparece como obligación legal en un código civil, quizá no sea tan religioso como pensábamos. Quizá sea, simplemente, humano. Quizá la biología, la ley y la fe estén diciendo lo mismo con distintos acentos: que el amor de verdad tiene una forma, y que esa forma incluye permanencia, exclusividad y apertura a la vida.

Esto no debería sorprendernos. En la primera parte de este libro vimos que nuestro cuerpo ya viene equipado para amar así: con mecanismos hormonales que nos empujan a vincularnos, a quedarnos, a cuidar.\footnote{Cf. capítulos 2 y 3 de este libro.} Si el cuerpo habla de permanencia, si la ley habla de permanencia y si la fe habla de permanencia\ldots{} quizá sea porque la permanencia no es un invento de nadie, sino algo que está en la naturaleza misma del amor.

Dicho esto, vamos al grano. Porque este capítulo va sobre la primera de esas tres columnas. La más temida. La más malentendida. Y, cuando se mira bien, la más hermosa.

Hablemos de la indisolubilidad.

\section{El vínculo que no se cuestiona}

Pero antes, pensemos un momento en algo que damos por sentado y que, curiosamente, nadie discute.

De todos los vínculos familiares que tienes, absolutamente todos son indisolubles. Haga lo que haga tu hijo, seguirá siendo tu hijo. Siempre. Aunque os peleéis, aunque no os habléis durante años, aunque te decepcione de la peor manera posible: es tu hijo. Lo mismo con tu padre, con tu madre, con tus hermanos, con tus primos. No puede pasar nada ---nada--- que haga que tu primo deje de ser tu primo. No hay papeles que firmar, no hay juez que lo decrete, no hay ley que lo permita. Los vínculos de sangre son para siempre, y a nadie se le ocurre cuestionarlo.

Y ninguno de esos vínculos lo elegiste. Te vinieron dados. No decidiste nacer en esa familia, no escogiste a tus padres, no firmaste ningún contrato con tus hermanos. Simplemente, así son las cosas.

Y sin embargo ---aquí viene la paradoja---, pensamos que podemos deshacer el único vínculo familiar que decidimos crear libremente. El único que elegimos con plena conciencia, con libertad, delante de testigos. El único para el que, como hemos visto en la primera parte de este libro, estamos biológicamente diseñados. El que garantiza la supervivencia de nuestras crías y, como veremos, nuestra propia santidad.

Los vínculos que no elegimos son indisolubles sin discusión. El vínculo que elegimos libremente, ese sí creemos que podemos deshacerlo. ¿No es extraño?

Quizá la pregunta no sea si la indisolubilidad es razonable. Quizá la pregunta sea por qué nos parece razonable la disolubilidad.

\section{La cadena}

Hay pocas palabras que generen más rechazo inmediato que «indisolubilidad». Suena a cerrojo. A candado sin llave. A puerta que se cierra de golpe y ya no se abre jamás.

La reacción es comprensible. Vivimos en una cultura que celebra la libertad de elegir y, sobre todo, la libertad de dejar de elegir. Podemos cambiar de trabajo, de ciudad, de teléfono, de look, de pareja. La posibilidad de salir es lo que nos hace sentirnos libres. ¿Qué clase de libertad es la que te dice «esto es para siempre, sin marcha atrás»?

Y la objeción más frecuente, la que todo el mundo lleva en la punta de la lengua, es esta: «¿Y si no funciona?»

Es una pregunta razonable. Todos conocemos matrimonios rotos. Todos hemos visto el dolor de una separación, la amargura de quien se siente atrapado, el sufrimiento de los hijos. Ante eso, la indisolubilidad parece no solo anticuada, sino cruel. ¿Cómo puedes exigir a alguien que permanezca en algo que le hace daño?

Entiendo la objeción. De verdad. No la despacho con un gesto de la mano. El sufrimiento en un matrimonio roto es real, y merece todo el respeto y toda la compasión del mundo.\footnote{El Papa Francisco dedica todo el capítulo VIII de \textit{Amoris Laetitia} a acompañar la fragilidad, reconociendo las «situaciones imperfectas» y pidiendo a la Iglesia que sea «hospital de campaña» más que «aduana»: «Se trata de integrar a todos, se debe ayudar a cada uno a encontrar su propia manera de participar en la comunidad eclesial.» Francisco, Exhortación apostólica \textit{Amoris Laetitia}, n. 297 (2016).}

Pero quiero proponerte algo. Quiero pedirte que dejes en suspenso, solo por unos minutos, esa objeción. Que no la olvides ---la recogeremos después---, pero que antes mires la indisolubilidad desde otro ángulo. Uno que quizá no hayas considerado.

\section{¿Y si fuera al revés?}

¿Y si la indisolubilidad no fuera una cadena que te ata?

¿Y si fuera un regalo que te dan?

Piénsalo así. Imagina que alguien te dice: «Yo me quedo contigo. Pase lo que pase. En lo bueno y en lo malo. Cuando estés guapo y cuando no lo estés. Cuando seas encantador y cuando seas insoportable. Cuando todo vaya bien y cuando todo se tuerza. No me voy. No te voy a dejar. Nunca.»

¿Quién no querría escuchar eso?

No estoy hablando de una fantasía romántica. Estoy hablando de la necesidad más profunda del ser humano: la seguridad de ser amado sin condiciones. De saber que hay alguien en el mundo que ha decidido no irse. Que ha puesto su libertad ---lo más valioso que tiene--- al servicio de una promesa: «Contigo, hasta el final.»

Eso es la indisolubilidad. No es lo que la Iglesia te exige. Es lo que tu cónyuge te regala.

\section{La promesa que el mundo no sabe hacer}

Vivimos rodeados de promesas con fecha de caducidad. «Te quiero» dura mientras las cosas vayan bien. «Estoy contigo» lleva una letra pequeña que dice «salvo que encuentre algo mejor». Las relaciones se viven con una especie de periodo de prueba permanente: si funciona, genial; si no, cada uno por su lado.

Y parece muy razonable. Muy maduro. Muy libre.

Pero ¿lo has pensado bien? ¿Cómo vives cuando sabes que la otra persona puede irse en cualquier momento?

Vives auditando. Vigilando. Midiendo. ¿Le sigo gustando? ¿Sigo siendo suficiente? ¿Estará pensando en otro, en otra? ¿He engordado demasiado? ¿He dicho algo que no debía? ¿Estoy a la altura?

Vives, en el fondo, con miedo.

Y el miedo es lo contrario de la libertad.

Ahora imagina lo otro. Imagina que esa persona te ha dicho «no me voy», y lo ha dicho de verdad, y tú le crees. Que sabes, con la certeza que da una promesa pública y solemne, que mañana va a seguir ahí. Y pasado mañana. Y dentro de diez años. Y cuando tengas arrugas y el pelo blanco y ya no te acuerdes de dónde has puesto las llaves.

¿Qué pasa entonces?

Entonces puedes quitarte la máscara.

\section{La libertad de ser tú}

Solo en la seguridad de que el otro no se va puedes mostrarte como eres de verdad.

Piénsalo. En las primeras citas ---todos lo hemos hecho--- te presentas en tu mejor versión. Te peinas, te arreglas, mides tus palabras, controlas tus manías. Escondes lo que crees que no gusta. Muestras lo que crees que impresiona. Estás actuando. No mintiendo, exactamente, pero sí editando. Poniendo filtro.

Y tiene sentido. Cuando la otra persona aún no se ha comprometido, cuando todavía está decidiendo si se queda o se va, tú no puedes permitirte el lujo de ser vulnerable. Mostrar tus debilidades es un riesgo: ¿y si al verlas decide que no mereces la pena?

Ahora imagina que la decisión ya está tomada. Que hay una promesa. Que esa persona ha dicho, delante de su familia, de sus amigos y de Dios: «No me voy.»

Entonces puedes dejar de actuar. Puedes decir «estoy agotado y hoy no puedo con nada» sin miedo a que te dejen. Puedes confesar «me equivoqué, perdóname» sin miedo a que sea el fin. Puedes mostrar tus cicatrices, tus inseguridades, tus rincones más oscuros, y saber que la otra persona no va a salir corriendo. Porque ya decidió quedarse.

Eso no es una cadena. Eso es el hogar.

Hay una diferencia enorme entre una relación en la que estás siempre haciendo casting ---buscando que te renueven el contrato--- y una en la que sabes que tienes tu sitio. En la primera vives tenso, vigilante, midiendo cada paso. En la segunda puedes respirar. Puedes crecer. Puedes ser tú.

Y lo paradójico es que solo cuando eres tú de verdad ---sin filtros, sin máscaras, sin poses--- el otro puede amarte de verdad. Porque ama lo que eres, no lo que finges ser.

La indisolubilidad no limita la libertad. La hace posible.

\section{La biología ya lo sabía}

Si has leído la primera parte de este libro, quizá estés pensando: «Esto me suena.» Y con razón.

Cuando hablamos de la química del vínculo,\footnote{Cf. capítulo 3 de este libro.} vimos que el enamoramiento ---esa tormenta de dopamina, noradrenalina y serotonina baja--- tiene fecha de caducidad. Dura entre uno y tres años, aproximadamente. Pasado ese tiempo, el cóctel cambia. Y lo que aparece no es vacío, sino algo distinto y más profundo: el apego.

El apego se construye con oxitocina y vasopresina, las hormonas de la confianza, la calma y la permanencia. Son las mismas hormonas que vinculan a una madre con su bebé. Y hacen algo que la dopamina del enamoramiento no puede hacer: crear un vínculo duradero, sereno, basado no en la excitación sino en la seguridad.

La biología, literalmente, nos empuja de la pasión al apego. Del «no puedo vivir sin ti» al «quiero construir contigo». Del terremoto a la tierra firme.

¿No es curioso? Nuestro propio cuerpo está diseñado para la permanencia. La evolución ha invertido millones de años en crear un sistema que nos lleva, paso a paso, del deseo inicial al vínculo estable. Porque nuestras crías ---esas criaturas indefensas que necesitan años de cuidado--- necesitan padres que se queden. No padres que se enamoren un rato y luego se vayan.

La indisolubilidad no es una ocurrencia de la Iglesia. Es la palabra que la fe le pone a algo que el cuerpo ya sabe.

\section{El regalo que transforma}

Hay algo más en la indisolubilidad, algo que es fácil de pasar por alto y que a mí me parece decisivo.

Cuando alguien te promete que no te va a dejar nunca, esa promesa te cambia. No en abstracto, no en teoría. Te cambia de verdad. Porque al recibirla, nace en ti algo que antes no estaba: el deseo de merecerla.

Nadie te lo pide. Nadie te lo exige. Pero cuando alguien ha puesto su vida entera en tus manos ---cuando ha apostado todo lo que tiene a una carta, y esa carta eres tú---, algo dentro de ti se moviliza. Quieres estar a la altura. No por obligación, sino por amor. No porque te castiguen si fallas, sino porque no quieres fallar a quien te ha dado tanto.

«Me ha prometido que se queda. Pues yo quiero ser alguien que merezca que se quede.»

¿Ves cómo funciona? La indisolubilidad no es un mandato que te aplasta. Es un regalo que te eleva. Te convierte en mejor persona, no porque te obligue, sino porque te inspira. Como el que recibe una confianza enorme y, precisamente por eso, se esfuerza en no defraudarla.

Es lo que ocurre con todas las grandes promesas. El padre que le dice a su hijo «siempre estaré aquí» no se está poniendo una cadena: se está dando una razón para ser fuerte. El amigo que dice «cuenta conmigo pase lo que pase» no está firmando una condena: está eligiendo ser mejor persona. La promesa de permanencia no te encadena a la miseria. Te empuja hacia arriba.

\section{Pero ¿y si no funciona?}

Lo prometí: íbamos a volver a esta pregunta. Porque no sería honesto esquivarla.

No todo matrimonio es un camino de rosas. Hay matrimonios donde hay dolor real, incluso violencia. Hay situaciones de las que hay que salir, y la Iglesia lo sabe. No estoy diciendo que haya que aguantar cualquier cosa en nombre de la indisolubilidad. Eso sería una caricatura ---y una crueldad---.

Pero fijaos en algo que pasa constantemente y que dice más que cualquier argumento teórico. Cuando una pareja se separa ---da igual si fue de mutuo acuerdo o entre abogados, da igual si se llevan bien o mal---, escuchadlos hablar. Prestad atención a cómo se refieren al otro. Es llamativo: muy a menudo, incluso años después de la separación, siguen diciendo «mi marido», «mi mujer». No «mi ex». No «la persona con la que estuve». \emph{Mi marido. Mi mujer.} Lo dicen con cariño o con rabia, con nostalgia o con amargura, pero lo dicen. Algo en ellos ---más fuerte que la distancia, más fuerte que el conflicto, más fuerte que los papeles del juzgado--- sigue reconociendo un vínculo que no han podido romper del todo. Como si el lenguaje supiera lo que la razón a veces no quiere admitir: que lo que hubo entre ellos no fue un contrato que se rescindió, sino un lazo que, de alguna manera misteriosa, sigue ahí. El vínculo es más fuerte de lo que pensamos. Más fuerte, incluso, de lo que quienes lo viven están dispuestos a reconocer.

Pero sería deshonesto dejar las cosas aquí, en la anécdota entrañable del lenguaje que delata un vínculo. Porque la realidad del fracaso matrimonial es más compleja y, a veces, más oscura.

Hay matrimonios que se rompen y que dejan heridas profundas. Hay personas que han vivido la indisolubilidad no como un regalo sino como una losa, como una condena a permanecer en algo que les hacía daño. Hay separaciones que no son fracasos sino actos de dignidad: la dignidad de quien dice «hasta aquí» ante la violencia, el desprecio o la destrucción sistemática de la persona. Y hay que tener la honestidad de nombrarlo.

La Iglesia, que a veces parece hablar solo del ideal, sabe esto perfectamente. La separación física está prevista en el derecho canónico precisamente para proteger a quien sufre.\footnote{CIC, can. 1153 §1: «Si uno de los cónyuges pone en grave peligro espiritual o corporal al otro o a la prole, o de otro modo hace demasiado dura la vida en común, proporciona al otro un motivo legítimo para separarse.»} Y existe un proceso de nulidad que no es un divorcio católico encubierto, sino algo muy distinto: la constatación, tras un examen riguroso, de que lo que parecía un matrimonio no llegó a serlo nunca ---porque faltó libertad, madurez suficiente, o un consentimiento verdadero---.\footnote{Cf. CIC, cann. 1095--1107. La declaración de nulidad no dice «este matrimonio ha terminado», sino «este matrimonio nunca existió válidamente». Es una distinción sutil pero fundamental.} No se trata de buscar una puerta trasera, sino de reconocer una verdad.

Y hay algo más que conviene decir, porque se dice poco: el fracaso de un matrimonio no es el final de la historia de amor de esa persona. Como en cualquier otra dimensión de la vida ---la profesional, la vocacional, la personal---, el fracaso puede ser una herida, pero también un punto de partida. Hay personas separadas que han encontrado, a través de ese dolor, una profundidad humana y espiritual que no tenían antes. El Papa Francisco lo dice con una claridad que la Iglesia no siempre ha tenido: esas personas «no están excomulgadas», y la comunidad cristiana debe ser para ellas «hospital de campaña» y no «aduana».\footnote{Francisco, \textit{Amoris Laetitia}, nn. 243, 291, 296--297 (2016). Todo el capítulo VIII de la exhortación está dedicado a «acompañar, discernir e integrar la fragilidad».}

Nada de esto debilita la indisolubilidad. La hace más real, más humana, más honesta. Porque la promesa «para siempre» no se pronuncia en un mundo perfecto sino en este, donde las personas son frágiles, cometen errores y a veces se equivocan incluso en lo más importante. Reconocer esa fragilidad no es renunciar al ideal; es amarlo con los ojos abiertos.

Lo que estoy diciendo es otra cosa. Estoy diciendo que la pregunta «¿y si no funciona?», cuando se hace \emph{antes} de casarse, como miedo preventivo, como excusa para no comprometerse, revela algo importante: revela que estamos mirando la indisolubilidad desde el lado equivocado.

La estamos mirando como una cláusula contractual que te atrapa. Pero no es eso. Es una declaración de amor que te sostiene.

¿Y si, en lugar de preguntar «qué pasa si no funciona», preguntáramos «qué pasa si sí funciona»? ¿Qué pasa si alguien de verdad cumple esa promesa? ¿Qué pasa si tienes la experiencia de ser amado así, sin reservas, sin salida de emergencia, sin condiciones?

Te diré lo que pasa: que tu vida se transforma. Que por primera vez puedes bajar la guardia del todo. Que descubres una libertad que no habías conocido: la libertad de ser amado no por lo que haces, sino por lo que eres. La libertad de no tener que ganarte cada día el derecho a ser querido.

Y esa experiencia ---la de ser amado incondicionalmente por alguien que ha elegido quedarse--- es quizá la imagen más cercana que tenemos del amor de Dios.\footnote{«El amor conyugal comporta una totalidad en la que entran todos los elementos de la persona ---reclamo del cuerpo y del instinto, fuerza del sentimiento y de la afectividad, aspiración del espíritu y de la voluntad---; mira a una unidad profundamente personal que, más allá de la unión en una sola carne, conduce a no tener más que un corazón y un alma.» Juan Pablo II, \textit{Familiaris Consortio}, n. 13.}

\section{Una promesa que viene de más arriba}

Quiero terminar con algo que quizá ya intuyes.

Si la indisolubilidad es un regalo tan grande, ¿de dónde viene? ¿Es solo una decisión humana, un acto de voluntad heroica? ¿O hay algo más?

Cuando Jesús habla del matrimonio, no inventa nada nuevo. Hace algo mucho más audaz: vuelve al principio. Le preguntan por el divorcio, y en lugar de discutir la ley de Moisés, remite a sus oyentes al Génesis: «Al principio no fue así.»\footnote{Mt 19,8. Los fariseos le preguntan por qué Moisés permitió el divorcio, y Jesús responde: «Por la dureza de vuestro corazón os permitió Moisés repudiar a vuestras mujeres; pero al principio no fue así.»} Al principio, cuando Dios pensó el amor entre un hombre y una mujer, lo pensó así: para siempre. No como castigo. No como prueba. Como regalo.

«Lo que Dios ha unido, que no lo separe el hombre.»\footnote{Mt 19,6.}

Fíjate: no dice «lo que la Iglesia ha unido». Dice «lo que Dios ha unido». El vínculo no lo crea la ceremonia, ni el sacerdote, ni el registro civil. El vínculo viene de arriba. Y precisamente por eso tiene una solidez que ninguna promesa puramente humana podría tener.

Esto es lo que la tradición cristiana llama el vínculo sacramental. Cuando dos bautizados se casan, su unión no solo es un compromiso entre ellos: es un signo del amor de Cristo por su Iglesia.\footnote{Cf. Ef 5,31-32. El Concilio de Trento definió la sacramentalidad del matrimonio y su indisolubilidad intrínseca. Cf. Concilio de Trento, Sesión XXIV (1563), Decreto \textit{Tametsi}. El \textit{Catecismo de la Iglesia Católica}, n. 1640, dice: «El vínculo matrimonial es establecido por Dios mismo, de modo que el matrimonio celebrado y consumado entre bautizados no puede ser disuelto jamás.»} Un amor que no tiene fecha de caducidad. Un amor que no se retira cuando el otro falla. Un amor que permanece.

Y la buena noticia es esta: no estáis solos en esa promesa. No se trata de que vosotros, con vuestra sola fuerza de voluntad, tengáis que sostener un compromiso de por vida. Se trata de que Dios mismo está en ese vínculo, sosteniéndolo desde dentro. La gracia del sacramento no es un adorno: es una fuerza real que trabaja en vuestro matrimonio, que os empuja a amar cuando ya no os quedan fuerzas, que os sostiene cuando todo parece venirse abajo.\footnote{«La gracia propia del sacramento del Matrimonio está destinada a perfeccionar el amor de los cónyuges, a fortalecer su unidad indisoluble. Por medio de esta gracia "se ayudan mutuamente a santificarse".» \textit{Catecismo de la Iglesia Católica}, n. 1641.}

No os piden que seáis héroes. Os ofrecen una mano que os sostiene.

\section{La promesa más hermosa}

Déjame resumirlo con una imagen.

Imagina dos formas de cruzar un río. En la primera, cruzas sobre piedras sueltas. Cada paso es incierto. No sabes si la próxima piedra aguantará tu peso. Avanzas con miedo, midiendo, calculando. Si la piedra se mueve, te caes.

En la segunda, cruzas por un puente de piedra. Sólido. Firme. Construido para durar. No tienes que pensar en cada paso. Puedes mirar alrededor, disfrutar del paisaje, caminar con confianza. El puente aguanta.

La indisolubilidad es el puente.

No te quita la libertad de caminar. Te da una superficie firme sobre la que hacerlo. No te obliga a llegar al otro lado: te permite llegar.

Y lo más hermoso de todo es que ese puente no te lo construyes tú solo. Lo construís juntos. Piedra a piedra. Día a día. Con cada «me quedo» silencioso, con cada perdón concedido, con cada mañana en la que decides amar aunque no tengas ganas. Y sobre ese puente, sostenido por una gracia que viene de más arriba, podéis caminar seguros. Juntos. Hasta el final.

Esa es la promesa que el mundo no sabe hacer. La promesa que suena a cadena y resulta ser alas. La promesa que nadie te va a hacer\ldots{} excepto la persona que de verdad te ama.

\paracaminar

\begin{enumerate}
  \item ¿Cómo percibíais la indisolubilidad antes de leer este capítulo: como una carga o como un regalo? ¿Ha cambiado algo en vuestra forma de verla? Hablad con sinceridad.

  \item Pensad en vuestra relación: ¿hay momentos en los que os habéis «puesto la máscara» por miedo a que el otro no aceptara lo que veía? ¿Qué pasaría si pudierais quitárosla del todo, con la seguridad de que el otro no se va?

  \item Intentad deciros, mirándoos a los ojos, la promesa de este capítulo con vuestras propias palabras. No tiene que ser solemne ni perfecto. Solo verdadero.
\end{enumerate}

\vinetacapitulo{ch08}
